\subsection{The primal toric-code tensor in the parity basis}

\begin{definition}[Primal toric-code map and virtual parity]%
    \label{def:peps_toric_code_primal_virtual_symmetry}
    \lean{TNLean.PEPS.toricCodeZ,
        TNLean.PEPS.toricCodePrimalParityMatrix,
        TNLean.PEPS.toricCodePrimalVirtualRep,
        TNLean.PEPS.toricCodePrimalVirtualRep_generator,
        TNLean.PEPS.mem_invariants_toricCodePrimalVirtualRep_iff,
        TNLean.PEPS.toricCodeEvenProjection,
        TNLean.PEPS.toricCodeEvenProjection_eq_half_one_add,
        TNLean.PEPS.toricCodePrimalMatrix}
    \leanok
    \uses{def:peps_quantum_double_primal}
    Write $G=\mathbb Z_2$ additively and put $V=G^4$.
    The primal tensor of
    \cite[Appendix~A, ``The Toric Code and quantum double models'']{Cirac2021Matrix}
    defines the map $P:\C^V\to\C^V$ with entries
    \begin{align}
        P_{s,c}=\delta_{c,f(s)},\qquad
        f(a,b,c,d)=(a+b,b+c,d+c,a+d).
    \end{align}
    In the basis $\ket0,\ket1$, let $Z=\operatorname{diag}(1,-1)$.
    The virtual representation of $G$ has generator
    $J=Z^{\otimes4}$. Its invariant subspace is
    \begin{align}
        V_{\mathrm{even}}
        =\{x\in\C^V:x_{t,r,b,l}=0\text{ whenever }t+r\ne l+b\}.
    \end{align}
    The orthogonal projection onto this subspace is
    \begin{align}
        \Pi_{\mathrm{even}}
        =\operatorname{diag}\bigl(\mathbf1_{t+r=l+b}\bigr)
        =\tfrac12(I+J).
    \end{align}
\end{definition}

\begin{theorem}[Isometry of the primal toric-code tensor]%
    \label{thm:peps_toric_code_primal_g_isometric}
    \lean{TNLean.PEPS.toricCodePrimalMatrix_conjTranspose_mul_self,
        TNLean.PEPS.isGIsometric_toricCodePrimalTensor}
    \leanok
    \uses{def:peps_toric_code_primal_virtual_symmetry, def:peps_g_injective}
    The primal toric-code map satisfies
    \begin{align}
        P^\dagger P=2\Pi_{\mathrm{even}},\qquad PJ=P.
    \end{align}
    Consequently it is $G$-injective and $G$-isometric for the virtual
    action generated by $Z^{\otimes4}$, with the precise normalization
    \begin{align}
        \langle Px,Py\rangle=2\langle x,y\rangle
        \qquad(x,y\in V_{\mathrm{even}}).
    \end{align}
    The factor $2$ records the unnormalized entries of the source tensor
    \cite{gap:rmp_peps_quantum_double_g_isometry}.
\end{theorem}

\begin{proof}\leanok
    \uses{thm:peps_quantum_double_primal_fusion}
    For a fixed virtual label $(t,r,b,l)$, the equation
    $f(s)=(t,r,b,l)$ has no solution unless $t+r=l+b$.
    Under this condition its solutions are exactly
    \begin{align}
        s=(t+h,h,h+r,l+t+h),\qquad h\in G.
    \end{align}
    There are therefore two solutions. Since $P_{s,c}=\delta_{c,f(s)}$,
    \begin{align}
        (P^\dagger P)_{c,c'}
        =\sum_s\delta_{c,f(s)}\delta_{c',f(s)}
        =2\delta_{c,c'}\mathbf1_{t+r=l+b}.
    \end{align}
    All nonzero entries have even virtual parity, so $PJ=P$.
    If $x\in V_{\mathrm{even}}$, then $P^\dagger Px=2x$;
    hence $Px=0$ implies $x=0$. For invariant $x,y$, the same identity gives
    $\langle Px,Py\rangle=\langle x,P^\dagger Py\rangle
        =2\langle x,y\rangle$.
\end{proof}
